\documentclass[11pt]{article}

% --- Packages ---
\usepackage[utf8]{inputenc}
\usepackage[margin=1in]{geometry}
\usepackage{enumitem}
\usepackage{booktabs}
\usepackage{titlesec}
\usepackage{hyperref}

% --- Custom Header ---
\titleformat{\section}{\large\bfseries}{}{0em}{}[\titlerule]

\begin{document}

\begin{center}
    {\huge \textbf{Assignment 1: AI System Specification}} \\
    \vspace{0.2cm}
    {\large \textbf{Course:} Introduction to Artificial Intelligence} \\
    \textbf{Topic:} PEAS Framework and Task Environment Analysis \\
    \vspace{0.5cm}
\end{center}

\section{Objectives}
The goal of this assignment is to apply the formal frameworks from \textit{Russell \& Norvig (Chapter 2)} to a real-world problem. You will demonstrate your ability to specify success metrics and classify the complexity of an AI's operational environment.

\section{Task Description}
Select \textbf{one} of the following modern AI applications (or propose one for approval):
\begin{itemize}[noitemsep]
    \item \textbf{A:} An Automated Warehouse Sorting Robot (e.g., Amazon Robotics).
    \item \textbf{B:} A Personalized News Feed Recommender (e.g., Twitter/X or TikTok algorithm).
    \item \textbf{C:} An Agricultural Monitoring Drone for crop disease detection.
\end{itemize}

\vspace{0.3cm}
\noindent For your chosen system, complete the following three tasks:

\section{Task 1: PEAS Specification}
Define the \textbf{PEAS} framework for your agent. Ensure your descriptions are specific (e.g., do not just say "Cameras," specify "High-resolution RGB and Infrared sensors").

\begin{description}
    \item[Performance Measure:] How exactly is success measured? List at least four metrics.
    \item[Environment:] Where does the agent act? What are the key external factors?
    \item[Actuators:] What physical or digital tools does the agent use to change its world?
    \item[Sensors:] What are the inputs the agent uses to perceive its state?
\end{description}

\section{Task 2: Environment Classification}
Classify the task environment according to the following dimensions. For each, provide a \textbf{one-sentence justification} explaining your choice.
\begin{itemize}
    \item Fully Observable vs. Partially Observable
    \item Single-agent vs. Multi-agent
    \item Deterministic vs. Stochastic
    \item Episodic vs. Sequential
    \item Static vs. Dynamic
    \item Discrete vs. Continuous
\end{itemize}

\section{Task 3: Critical Analysis}
Answer the following questions in short-paragraph format:
\begin{enumerate}
    \item Which of the six dimensions mentioned in Task 2 poses the \textbf{greatest technical challenge} for an AI engineer building this system? Why?
    \item Propose one \textbf{Utility Function} that would help the agent handle a trade-off (e.g., Speed vs. Safety). How would the agent's behavior change if the weight of this utility was doubled?
\end{enumerate}

\section{Submission Guidelines}
\begin{itemize}
    \item Submit your response as a single \textbf{PDF} file via email.
    \item Format: 12pt font, maximum 2 pages.
    \item Academic Integrity: All justifications must be your own. Use of AI for formatting is allowed, but the design logic must be original.
\end{itemize}

\end{document}